\section{Counting the cuspidal spectrum that receives the deformation points}
\label{sec:spectrum-deformation}

For a CM number field $K$, set $d_K=[K:\mathbb Q]$. Let $n\ge2$,
$M\ge1$, $H\ge0$ be integers, and let $m_0>\cdots>m_{n-1}$ be
integers satisfying $|m_i|\le H$. Write $K_n(M)$ for the principal
congruence subgroup of $\mathrm{GL}_n(\widehat{\mathcal O}_K)$ of level
$M$, viewed inside the finite adelic group. The collection
$\mathrm{Sig}_K(m,M)$ consists of cuspidal representations with
contragredient equal to the complex-conjugation pullback in the sense of
the source, a nonzero $K_n(M)$-fixed vector, and infinitesimal character
\[
                    I_m=\{(n-1)/2+m_i:0\le i<n\}
\]
at every embedding $K\hookrightarrow\mathbb C$. For fixed $K,n$
and any finite set $S$ of pairwise nonisomorphic members, we will obtain
\begin{equation}
 |S|\le C(K,n) M^{s}(H+1)^s,\qquad s=d_K n^2 .
 \label{eq:spectrum-count}
\end{equation}
Here the bounds in the source construction are kept separate until
the final multiplication. The quantified source theorem
\texttt{card\_k\_of\_printed\_results} permits unnamed constants and
uses the coarser exponent $s+s$ on $M(H+1)$. Both estimates follow from
the same two inequalities. The constants depend on $K,n$.

\subsection{One coefficient system for all the representations in $S$}

For $0\le k<n$ put $\nu_k=k+1+H-m_{n-1-k}$ and $H'=n+2H$.
These are nonnegative integers bounded by $H'$. They are
nonincreasing, since
\[
       \nu_{k+1}-\nu_k=1-(m_{n-2-k}-m_{n-1-k})\le0 .
\]
Let $X$ be a generic $n$ by $n$ matrix. For $0\le h<H'$ the set
$\{k:\nu_k>h\}$ is an initial segment of size $b_h$.
Take the product $f(X)$ of the top-left principal minors of orders
$b_h$, with the empty minor equal to one. This polynomial is nonzero
by evaluation at the identity, and has degree at most $nH'$.
For an upper triangular $B$, the top-left $b_h$ by $b_h$ block of
$XB$ is the corresponding block of $X$ multiplied by that of $B$.
Consequently
\[
             f(XB)=f(X)\prod_{i=0}^{n-1} B_{ii}^{\nu_i}.
\]
The space of polynomials of degree at most $nH'$ in the $n^2$
entries of $X$ is stable under the right substitution $X\mapsto XB$.
Each monomial has each exponent between zero and $nH'$, giving
dimension at most $(nH'+1)^{n^2}$. The highest-weight scalar
calculation in \texttt{Weyl.existsScalar} extracts a nonzero module $P$
of dimension at most this bound, with the central quantum determinant
acting as $\prod_k(u-k+\nu_k)$.

Use one copy at every embedding of $K$, with the determinant shift
$N_0=n+H$ in \texttt{exists\_tens}. Their external tensor product
has dimension at most $(nH'+1)^{n^2d_K}$. The logarithmic differential
at a given embedding acts through that factor. The central polynomial
therefore acts by the scalar
\[
 \prod_{k=0}^{n-1}(u-N_0-k+\nu_k)
       =\prod_{k=0}^{n-1}(u-(n-1)-m_{n-1-k}).
\]
An irreducible submodule $\Xi$ inherits this scalar, and the dimension
can only decrease. In the infinitesimal-character convention the
polynomial is evaluated at $u+(n-1)/2$. Its value is
$\prod_i(u-((n-1)/2+m_i))$, giving $\Xi$ the character $-I_m$
at every embedding. In particular the sign of the dual coefficient
system is fixed throughout $S$. Since
$n(n+2H)+1\le(n+1)^2(H+1)$, we have
\[
                 \dim\Xi\le((n+1)^2)^s (H+1)^s .
\]
This uses the highest-weight scalar calculation and tensor realization
in the source. The transformation law of $f$, dimension estimate,
reversal of indices and infinitesimal-character shift were computed
above.

\subsection{The level factor in cohomology}

Put $G=\mathrm{GL}_n(\mathbb A_{K,f})$ and
$\Gamma=\mathrm{GL}_n(K)$, acting on the left. For the fixed level
$U_0=K_n(3)$ there are finitely many double cosets
$\Gamma g_i U_0$. This is the finite-double-coset input, together
with the finite index of $U_0$ at integral level one. Write
$\Gamma_i=\Gamma\cap g_iU_0g_i^{-1}$. The arithmetic subgroup
cohomology theorem \texttt{bddArith} gives a constant $c_i$
such that for each finite-dimensional complex $\Gamma_i$-module $V$,
\[
          \sum_{j\ge0}\dim H^j(\Gamma_i,V)\le c_i\dim V .
\]
Its hypotheses apply because the $\Gamma_i$ are commensurable with
$\mathrm{GL}_n(\mathcal O_K)$ in $\Gamma$. The underlying geometric
input is a bounded arithmetic-group complex; the commensurability
comparison is part of that theorem.

Set $U'=K_n(3M)$. Every fiber of $G/U'\to G/U_0$ has
$D=[U_0:U']$ elements: the fiber above $gU_0$ is $gU_0/U'$.
For a coefficient representation of dimension $a$ the function
space on such a fiber has dimension $Da$. The inverse image of a
$\Gamma$-orbit in $G/U_0$ contributes a coinduced representation
from $\Gamma_i$, with this fiber function space as its coefficients.
Concretely a section is determined on each translate by its values
on the fiber at $g_iU_0$, subject to the $\Gamma_i$-action there.
Shapiro's lemma identifies its cohomology with that of $\Gamma_i$
on the fiber. Summing over the finitely many orbits gives
\[
 \sum_{j\ge0}\dim H^j\!\left(\Gamma,
        \operatorname{Map}(G/U',\mathbb C^a)\right)
                  \le \left(\sum_i c_i\right) Da .
\]
The action on functions is
$(\gamma f)(x)=\xi(\gamma)f(\gamma^{-1}x)$ for the coefficient
representation $\xi$. This is the cohomology quantity
\texttt{cohDim} used by the spectrum count.

Reduction modulo $3M$ embeds the quotient by its principal-congruence
kernel into $\mathrm{GL}_n(\widehat{\mathcal O}_K/
3M\widehat{\mathcal O}_K)$. The coefficient ring has
$(3M)^{d_K}$ elements by the rank of the integral lattice. The
number of invertible matrices is at most the number of all matrices,
so $D\le(3M)^s$. This proves
$\mathrm{cohDim}\le c_2 M^s a$ with
$c_2=3^s\sum_i c_i$, independent of the level and the coefficients.

Every $\pi\in S$ has the opposite infinitesimal character to $\Xi$.
The cohomological criterion \texttt{ClozelCohC} gives nonzero relative
Lie-algebra cohomology for $\pi\otimes\Xi$. Every fixed vector at
level $M$ remains fixed at level $3M$ because reduction modulo $M$
factors through reduction modulo $3M$. The cuspidal cohomology
injection \texttt{BorelInjC}, applied to the finite pairwise
nonisomorphic set $S$ and the open compact group $U'$, gives
$|S|\le\mathrm{cohDim}$ for $a=\dim\Xi$. Combining the inequalities
proves \eqref{eq:spectrum-count} with
$C(K,n)=c_2((n+1)^2)^s$. The argument applies to every finite set of
representatives and therefore also bounds the number of isomorphism
classes.

For rank three over the imaginary quadratic field used in the
single-character estimate, $s=18$. After the inverse cyclotomic
twist the weight vector is $(e-1,-1,-e-1)$, with $e\ge1$.
It decreases strictly, its opposite pairs sum to $-2$, and one
can take $H=e+1$. Thus the bound for this CM spectrum is at most
$C(K,3)2^{18}[M(e+1)]^{18}$, since $e+2\le2(e+1)$.
The count concerns global systems. Section~\ref{sec:common-trace-integer}
explains how global labels give traces at a fixed unramified prime.
The uniform passage over all residue characteristics also uses local
compatibility, Satake integrality and conjugation of the systems.

\subsection{Which mixed points contribute to the two spectra}

Write $\Omega=\overline{\mathbb Q}$, $G_F=\operatorname{Gal}(\Omega/F)$,
and use the deformation data of Section~\ref{sec:single-character}.
In this argument $K$ is the imaginary quadratic field cut out by
the quadratic multiplier $\eta$, so $\ker\eta=G_K$. The coefficient
ring $\mathcal O$ has absolute ramification index one here, and
$\iota:\overline{\mathbb Q}_\ell\simeq\mathbb C$ is fixed.
Let $r$ be a geometric point of the marked three-dimensional
representation and $X_g$ its matrix on $\ker\eta$. Index the rows
and columns by $0,1,2$. A mixed point is neither upper nor lower
triangular in the normalization. The theorem \texttt{twoPt\_shape}
gives two cases. In type I the $X_g$ act irreducibly on three
coordinates. In type II the middle coordinate and the outer pair
form blocks, and the two by two outer corner acts irreducibly.
The two cases cannot overlap because the middle coordinate is an
invariant line in type II.

Let $B$ be the chosen valuation ring of $\Omega$ above $\ell$.
The endomorphism and automorphy input \texttt{GdEndAlt} supplies one
marking with an integral finiteness conclusion and integral ordinary
endpoint statements for ranks three and two. The integral-point
theorem \texttt{exists\_ordPt\_of\_point} says precisely what is
needed to use these statements. For every permitted local deformation
condition, its marked ring, and an evaluation of the algebra of
matrix entries giving $r$, there exist a finite DVR
$\mathcal O\subset A\subset\overline{\mathbb Q}_\ell$, a unit $t$
of $\overline{\mathbb Q}_\ell$, and an ordinary $A$-point $r'$
whose scalar extension is the torus conjugate of $r$ by
$T=\operatorname{diag}(t,1,t^{-1})$. Its construction uses finiteness
of the weight-zero algebra, a graded finite-point theorem and the
local deformation conditions. The same marking works for every such
point. We spell out the use of conjugation because it changes the
coordinates of the outer corner.

On $\ker\eta$ one has $X'_g=TX_gT^{-1}$. In type I this carries
invariant subspaces bijectively to invariant subspaces. The ordinary
irreducibility theorem \texttt{matIrr\_open\_gd} upgrades the resulting
irreducibility to each open subgroup of $G_{\mathbb Q}$ contained in
$\ker\eta$. This stronger conclusion is an explicit hypothesis of
the rank-three ordinary endpoint. Applying it gives conjugate
self-dual automorphy for the restriction $\rho'=r'|_{G_K}$ after
the twist $\rho'\otimes\chi_\ell^{-1}$, including the source's
local witness for the continuous representation. Conjugating by
$T^{-1}$ intertwines this twist with
$\rho\otimes\chi_\ell^{-1}$. The same automorphy conclusion therefore
holds at the original point. Admission to the counted spectrum also
requires the level conditions and the local assertion denoted
$\mathrm{Unl}(\ell,M,(e,0,-e),\rho)$. The separate result
\texttt{admissible\_typeI\_even\_level\_alt} supplies these for the
admissible type-I data.

In type II, the point theorem gives $(X_g)_{11}=1$ for $g\in G_K$.
For an extension $W:G_{\mathbb Q}\to\mathrm{GL}_2(
\overline{\mathbb Q}_\ell)$ agreeing with the outer corner on $G_K$,
the marking requires $\operatorname{tr}W(c)=0$ at the designated
complex conjugation $c$. Let $D_t=\operatorname{diag}(t,t^{-1})$
and $W'=D_t W D_t^{-1}$. The odd-weight matrix entries of $X'_g$
still vanish, its middle entry is one, and its corner on $G_K$ is
$D_t\operatorname{corner}(X_g)D_t^{-1}=W'(g)$. Its corner remains
irreducible. Also $\operatorname{tr}W'(c)=0$. The rank-two ordinary
endpoint gives a classical newform $f$ of weight $2e+1$ attached to
$W'\otimes\chi_\ell^{-e}$, together with the witness
$\mathrm{WitAQ}(\iota,W')$ for weak admissibility. Conjugacy
transports both assertions to $W$. Thus the untwisted rank-three
restriction on $G_K$ has the form $\mathbf1\oplus W|_{G_K}$ and
the rational spectrum supplies this branch's labels. In the assembled
integer bound, the rank-two level is $M|\operatorname{disc}K|$ and
the integral weight vector is $(e,-e)$. The rank-three
CM count above applies to the type-I branch. The rational count
and its quadratic restriction argument appear in
Section~\ref{sec:common-trace-integer}.

\paragraph{Scope of the arithmetic inputs.}
The cohomological criterion, cuspidal cohomology injection, finite
double coset theorem, arithmetic-stabilizer bound, and scalar
highest-weight calculation enter the spectrum count. The passage from
integral ordinary points to automorphy uses the separate endpoint
theorem, proven in the release through linked ordinary systems,
descent, and unitary automorphy. The conjugacy argument above proves
the transfer to the original point. It uses the marked shape and
open-subgroup irreducibility theorems as inputs. The subsequent
rigidity argument applies adjoint Selmer vanishing to type I and
rank-two adjoint and Kato inputs to type II.

\paragraph{Source map.}
The definitions of the spectrum and weights are in
\texttt{AutomorphicForms.IsoOfFrob}. The coefficient construction is
\texttt{TWeylC}, with the minor polynomial in \texttt{TopPol} and the
tensor and scalar steps in \texttt{ExistsTens} and \texttt{ExistsScalar}.
The orbit argument is \texttt{TemperedLocalComponents.CohBoundFree}
and the final cohomological comparison is \texttt{TCardK}. The
rank-three specialization is in
\texttt{FrobeniusTraceCountBoundOfCmEigensystemCount}.
The point statements are \texttt{UnitaryThreeGroups.PtShape},
\texttt{CurveDVRSpecialization.WeightZeroAlgebraLeEntryAlgebra},
\texttt{MatIrrOpenGd}, \texttt{AutEnd3AltAlt}, and
\texttt{AutomorphyU3OfEndoscopy}. The precise modular-form
normalization is in \texttt{AutomorphyDescent.AutNewformW}.

\footnote{\href{https://github.com/Deicyde/riemann-hypothesis/blob/40e56bab026552e089d7afcfeff08d2defb9c005/src/Development/AutomorphicArithmetic/AutomorphicVarieties/ModularityArgument/ThresholdArguments/DenseThresholdCases/TCardK.lean}{The counting deduction},
\href{https://github.com/Deicyde/riemann-hypothesis/blob/40e56bab026552e089d7afcfeff08d2defb9c005/src/Development/AutomorphicArithmetic/AutomorphicVarieties/AutomorphicForms/TemperedLocalComponents/CohBoundFree.lean}{the level and cohomology exponents}, and
\href{https://github.com/Deicyde/riemann-hypothesis/blob/40e56bab026552e089d7afcfeff08d2defb9c005/src/Development/AutomorphicArithmetic/AutomorphicVarieties/ModularityArgument/ThresholdArguments/DenseThresholdCases/AutomorphyU3OfEndoscopy.lean}{the integral endpoint comparison}.}
